%% This section adds chapter 2
\section{Theoretical Background}
The following sections provide a theoretical background of the various technologies, protocols, standards and scientific approaches relevant to this thesis work.
\subsection{CBRN Defense Systems}
CBRN defense systems are systems which are designed to aid first responders with their strategic decision making, standard operating procedures and situational awareness before, during and after the occurance of CBRN incidents. CBRN defense refers to the defense against chemical, biological, radiological and nuclear threats. \par
CBRN defense systems are used by militaries, first responders and dedicated CBRN defense units. These systems are comissioned by governments in accordance with CBRN preparedness goals and strategies. Such systems typically consist of hardware and software components which work together to collect, process and present information related to CBRN incidents. The hardware can include CBRN detectors, sensors, communication devices, personal computers, laboratory equipment and various mechanical and electrical automation systems. The exact hardware depends on the purpose of the system and the environment it is designed to operate in. \citep{PolicastroA.J.2008_itP,alma991972120506254}
\par
The software components are used to collect and proccess information from the hardware components and provide it to the users of the system. A common function of these types of systems is to provide near real-time measurements and alerts from collected sensors and other information sources. Furthermore, the information is typically presented to users using a map-based interface, where measurements, detected hazards, personnel, equipment and other relevant information are shown in reference to their geographical location. This allows information from multiple sources to be presented to users in a way that improves their  situational awareness.
\par
CBRN defense systems typically record all available information within the system by design, which enables users to performpost event analysis of CBRN incidents. Reviewing and cross referencing plots of recorded measurements and locations, generating and managing reports, overall incident review and incident documentation are common activities. These features can be typically designed in accordance with organizational standard operating procedures and regulations, reducing the need for manual documentation and reporting and increasing the accuracy of results from such procedures. Furthermore, the ability to perform post analysis of CBRN incidents can be used to improve organizational performance and preparedness for future incidents.
\par
These systems need to provide information with low delay and sufficient accuracy for decision making. Reliability, promptness and availability are important features of these systems, as they may operate during situations in which the communication infrastructure or individual sub-components are unavailable or degraded. Each sub-component of the software system can introduce delays and inaccuracies, act as a bottleneck or a point of failure. Therefore, every layer of the system and interfaces between sub-components need to be designed with rigorous performance and reliability requirements.
\subsection{Low Earth Orbit (LEO) Satellite Networks}
Low Earth Orbit (LEO) satellites are satellites that orbit the Earth at a speed higher than that of Earth's rotation. As a result of that higher speed, LEO satellites appear to continuously circulate around the Earth. In order to provide global coverage a constellation of LEO satellites is required. \citep{Sun2014}
\par
The altitude of a satellite has a direct effect on the characteristics of the communication system. Compared to geostationary Earth orbit (GEO) satellites, LEO satellites operate considerably closer to the Earth. The shorter distance reduces the signal propagation time and the free-space path loss. This makes LEO satellites suitable for services where delay is important. The lower altitude also means that an individual satellite covers a smaller area and remains visible to a user for a limited time. Furthermore, a LEO network must use multiple satellites and coordinate the changing connectivity between satellites, user terminals and gateway stations. The resulting network is categorized as a mobile communication system rather than a single fixed satellite link of fixed infrastructure communication system. \citep{Sun2014}
\par
The following discussion of LEO satellite orbit mechanics, link characteristics and networking concepts draws primarily on the satellite networking principles presented by \citet{Sun2014}.

\subsubsection{Orbital and Constellation Characteristics}
An orbit is determined by the motion of a satellite under the influence of the Earth's gravity. The orbital altitude, inclination, eccentricity and position of the orbital plane determine the satellite's ground track and the areas that can be served. For communication systems, these parameters are selected together with the number and arrangement of satellites to provide the required coverage and availability. A single LEO satellite cannot provide continuous service over a large geographical area because its footprint moves across the Earths' surface. A constellation addresses this limitation by distributing satellites across several orbital planes so that another satellite can become available as the current satellite moves out of view. \citep{Sun2014}
\par
The movement of a LEO constellation is predictable, but the communication topology is dynamic. The relationships between satellites, gateways and user terminals change continuously as the satellites follow their orbits. In a constellation with inter-satellite links, the links within an orbital plane can remain relatively stable for a period of time, while links between orbital planes and links to ground stations change as the satellites move. This predictability can be used for network planning and routing, but it does not make the network static. Routing decisions and resource allocation must account for the expected lifetime and quality of available links. \citep{Sun2014}

\subsubsection{LEO Satellite Link Characteristics}
The communication path between an application and a remote endpoint may contain several segments: the user terminal to the satellite, a possible inter-satellite path, the satellite to a gateway, and the terrestrial connection to the destination. The propagation component of delay is generally lower for LEO systems than for GEO systems because the signal travels a shorter distance. However, the end-to-end delay also depends on gateway placement, routing, processing, scheduling and retransmissions. Therefore, the lower orbit should not be treated as a guarantee of a constant or minimal application-level latency. \citep{Sun2014}
\par
LEO links are also affected by the relative motion between the satellite and the user terminal. This motion changes the received signal conditions and produces Doppler shifts that must be compensated by the communication system. Atmospheric conditions, obstructions, antenna visibility, interference and variations in available radio resources can further affect the link. Link quality can consequently vary during an otherwise continuous application session. These variations are important for CBRN systems because a short period of degradation can delay an alert, interrupt a telemetry stream or trigger transport-layer loss recovery.
\par
Satellite networks traditionally use link-layer mechanisms such as error-correcting codes, retransmission and bandwidth allocation to improve the quality and availability of the radio link. These mechanisms can hide some transmission errors from higher layers, but they can also add delay and cause the observed round-trip time to vary. The transport protocol therefore operates over a path whose capacity, delay and loss characteristics may change over time. Measurements of an LEO application should differentiate between the radio-link behavior and the end-to-end behavior observed by the application. \citep{Sun2014}

\subsubsection{Mobility, Coverage and Handovers}
In a terrestrial fixed network, an endpoint can often remain attached to the same access point for the duration of a connection. This assumption does not generally hold for LEO networks. As a satellite moves across the sky, the user terminal may need to transfer its connection to another satellite or gateway in order to maintain service. A handover can involve establishing a new physical link, updating routing information and preserving the reachability of the user terminal. The network may also need to coordinate handovers between the space and ground segments. \citep{Sun2014}
\par
A handover does not necessarily cause a complete application outage, but it introduces a period during which the available path can change or packets can be delayed, reordered or lost. The practical effect depends on how the LEO provider implements mobility and on whether the user terminal retains the same network-layer identity. From the perspective of the transport layer, a path change can appear as a temporary increase in round-trip time, a reduction in throughput, packet loss or a connection interruption. These effects make handover behavior a relevant test condition for the resilience evaluation in this thesis.

\subsubsection{IP Networking and Transport Implications}
LEO satellite networks can be integrated with terrestrial IP networks through gateways and satellites may also contain routing functionality. In a constellation, routing is more complex than in a system consisting of one large-coverage satellite because the connectivity between nodes changes over time. Since satellite positions are predictable, orbit information can be used to anticipate topology changes and update routing decisions. Nevertheless, routing and mobility mechanisms must ensure that packets continue to reach the correct user terminal when its point of attachment changes. \citep{Sun2014}
\par
These characteristics create a distinction between link continuity and session continuity. A temporary disruption in one radio or network path does not necessarily mean that the communicating endpoints have become permanently unreachable. For a near-real-time CBRN application, the transport protocol should recover from short disruptions without requiring unnecessary application-level reinitialization, while still exposing enough information to manage stale or delayed data. This is one reason why connection establishment, loss recovery, multiplexing and behavior during path changes are important when selecting and evaluating a transport protocol over LEO satellite internet.
\par
The satellite network also introduces a shared and potentially variable communication resource. Competing traffic, scheduling decisions and gateway capacity can affect the service available to an individual user terminal. Quality of service is therefore not determined only by the nominal bandwidth of the satellite system. Latency, jitter, packet loss, service availability and the priority of individual data flows are also relevant. For CBRN telemetry, a small and time-sensitive alert may be more important than a larger non-critical payload, so the communication design should consider both the network conditions and the relative importance of the data.
\subsubsection{Starlink Satellite Network Characteristics}

\subsubsection{Challenges Relevant to CBRN Communication}
The main advantage of LEO satellite networks for CBRN defense is that they can provide connectivity in locations where terrestrial infrastructure is unavailable, damaged or deliberately disrupted. At the same time, the network has characteristics that complicate the delivery of near-real-time information. The most relevant challenges are variable delay, packet loss, changing throughput, intermittent visibility, handovers, gateway dependence and limited control over the underlying satellite network. These challenges can occur independently or at the same time, which makes average latency alone insufficient for evaluating the communication system.
\par
For this thesis, the relevant performance properties are end-to-end latency, connection establishment time, packet loss, recovery after temporary disconnection and the delivery behavior of high-priority telemetry during network changes. The LEO network is treated as an external communication environment whose detailed radio and constellation control mechanisms are not modified by the proposed artifact. The artifact instead operates at the transport and application boundaries, where it can measure the resulting network behavior and apply transport-level mechanisms to support resilient data delivery. The following section introduces QUIC as the transport protocol selected for this purpose.

\subsection{QUIC Protocol}


\clearpage  % This command will start the next section from the new page
